\section{Restoration of the Shared L9 Directions}
\label{sec:allstep-restoration}

Does undoing the perturbation along the shared directions move the commands back toward the clean run's? Restoration tests this at L9, writing clean values along these directions into the perturbed run at every generation step. It needs a clean observation of the same physical state, so it is a controlled test, not a deployable method. On held-out tasks, command error falls against the unedited run and both controls, and the shared directions reach most of what replacing the whole L9 state does. Task success is higher by 1.85 and 3.33 points (intervals in Table~\ref{tab:allstep-restoration-contrasts}); replacing the whole L9 state raises it further (post hoc, pointwise). This supports main Sec.~\ref*{M-sec:restoration}, Table~\ref*{M-tab:restoration} and Figs.~\ref*{M-fig:teaser}b and~\ref*{M-fig:native-restoration}.

Six runs per case answer three questions: does restoration help at all (against the unedited run), must the directions be the shared ones (against clean-feature and random controls), and how much can this layer do (replacing the whole L9 state)? Fixed in advance were the site, task success as the primary endpoint, command error as the second, and six simultaneous comparisons per endpoint. The controls write less; with edit size matched in development, restoration is better than clean-feature directions under Camera but not under Sensor Noise (Sec.~\ref{sec:dev-restoration-controls}).

\subsection{Directions, tasks and cases}

\paragraph{Directions and site.}
The edit acts on the action expert's L9 post-block state (ten ordered action positions, 1,024 channels). L9, the middle recorded layer, was fixed before the confirmation runs as a compromise between the response coverage of L1 and the command prediction of L18 (main Sec.~\ref*{M-sec:layers}); one site, and two of the four families (Camera and Sensor Noise), keep six runs per case tractable. The shared and clean-feature directions are the first sixteen directions of the position-centered LIBERO-Plus fits of Sec.~\ref{sec:mean-controls} (each plan position's mean removed before fitting), learned at the first query and first generation step and shared by all positions.

\paragraph{Held-out tasks.}
Development uses folds 0 and 3, confirmation folds 1, 2 and 4: nine and 31 tasks with no overlap, so that choices made in development are tested on unseen tasks. Each task uses the fit from its fold of the original split, which held out a task component (tasks sharing a clean scene) with a family; neither is in that fit's training data. Camera confirmation has 27 tasks and 48 settings; Sensor Noise has 30 tasks and 60 settings; the two share 26 tasks. Each setting has one start state, two policy-noise draws and six runs, giving $108\times2\times6=1{,}296$ episodes. All planned command measurements and outcomes are recorded.

\subsection{Restoration and the six runs}

\paragraph{A clean observation of each run's own current state.}
Runs move apart after the first query, so each needs a clean observation of its own state. Let $s_q^a$ be the physical and controller state of run $a$ at query $q$, and $x_{qj}^a,t_j$ its current plan and time at generation step $j$. Both observations are rendered from this $s_q^a$. For Camera, the clean observation restores the setting's complete clean visual configuration; for Sensor Noise, it removes the image transform from the same underlying image. Rendering it leaves the simulator, controller and random-number states unchanged.

Two forward passes with these observations and the same $x_{qj}^a,t_j$ give $H_T,H_C\in\mathbb R^{10\times1024}$ at L9, so $H_C$ is the clean value for the run's current state. Restoration sets
\begin{equation}
 D=H_T-H_C,\qquad
 H'=H_T-\alpha D B B^\top,\qquad \alpha=1,
 \label{eq:allstep-restoration}
\end{equation}
where $B\in\mathbb R^{1024\times16}$ has orthonormal columns. Up to the precision of the write, $H'B=H_CB$, and the orthogonal complement stays perturbed. Position centering is used only to fit $B$; the edit projects the full current difference $D$ onto the fixed directions. After the write, later blocks continue with the perturbed image prefix.

\paragraph{The six runs.}
All six runs start from the same simulator and controller state and use the same policy noise at each query. The clean run $N_C$ (the target) and the perturbed run $N_T$ are unedited. Restoration $R_{\mathrm{resp}}$ and the controls $R_{\mathrm{clean}}$ and $R_{\mathrm{rand}}$ apply Eq.~\eqref{eq:allstep-restoration} with the shared, clean-feature and random directions, at all ten generation steps of every policy query, with $\alpha=1$. $R_{\mathrm{full}}$ replaces the whole L9 state, $H'=H_C$, at the same site and schedule; its later blocks also read the perturbed prefix. Once an edit changes generation, each run follows its own plan and trajectory.

The random control uses a fixed mixture of three random sets of 16 orthonormal directions, assigned within task and family, with both noise draws of a setting receiving the same set. One set is the first sixteen columns of the 32-direction random set of Sec.~\ref{sec:mean-controls}; the other two come from the same generator with separately fixed seeds.

\begin{table*}[t]
\caption{\textbf{Restoration against the unedited run and both controls: command error falls in all six comparisons; task success is higher against the unedited run.} Restoration $R_{\mathrm{resp}}$ against each comparator. Positive values favor restoration: first five commands closer to those of the clean run $N_C$ (in units of a fixed family scale $S_f$), or more successes. Each endpoint has its own family of six simultaneous 95\% whole-task bootstrap intervals. Camera: 27 tasks, 48 settings and 96 cases; Sensor Noise: 30 tasks, 60 settings and 120 cases. Every success interval includes zero; against $R_{\mathrm{clean}}$ under Sensor Noise its lower end is zero.}

\centering\small
\setlength{\tabcolsep}{7pt}
\begin{tabular}{@{}llrr@{}}
\toprule
Family & Comparator & Command-error reduction $G_{f,b}$ [simultaneous interval] & Success change (pp) [simultaneous interval]\\
\midrule
Camera & $N_{\mathrm{T}}$ & $0.3068\;[0.1573,0.4821]$ & $+1.85\;[-8.33,11.11]$\\
Camera & $R_{\mathrm{clean}}$ & $0.2908\;[0.1423,0.4641]$ & $-1.85\;[-12.04,7.41]$\\
Camera & $R_{\mathrm{rand}}$ & $0.2999\;[0.1530,0.4724]$ & $+1.85\;[-9.26,12.96]$\\
\midrule
Sensor Noise & $N_{\mathrm{T}}$ & $0.0523\;[0.0179,0.0985]$ & $+3.33\;[-1.67,9.17]$\\
Sensor Noise & $R_{\mathrm{clean}}$ & $0.0337\;[0.0035,0.0740]$ & $+3.33\;[0.00,8.33]$\\
Sensor Noise & $R_{\mathrm{rand}}$ & $0.0516\;[0.0165,0.0996]$ & $+1.67\;[-2.50,6.67]$\\
\bottomrule
\end{tabular}

\label{tab:allstep-restoration-contrasts}
\end{table*}


\subsection{Endpoints and simultaneous intervals}

\paragraph{Command error.}
Let $A_{tvn}^a\in\mathbb R^{5\times3}$ hold the translation parts of the first five commands of the first query, for task $t$, setting $v$, noise draw $n$ and run $a$. The command error is
\begin{equation}
 e_{tvn}^a=\left\|A_{tvn}^a-A_{tvn}^{N_C}\right\|_F,
 \label{eq:allstep-command-error}
\end{equation}
the Euclidean norm over fifteen numbers, in command units after the adapter. All 1,296 records contain the five commands, logged after the environment step and identical to the first five positions of the ten-position plans. The clean target is paired by setting, start state and noise draw.

\paragraph{Task-equal contrasts and fixed scales.}
Positive contrasts favor restoration. Write $\langle z\rangle_f$ for the average over the two noise draws, then over settings within a task, then equally over the tasks of family $f$. For comparator $b\in\{N_T,R_{\mathrm{clean}},R_{\mathrm{rand}}\}$, the command-error reduction and the success change are
\begin{align}
 G_{f,b}&=\frac{\langle e^b-e^{R_{\mathrm{resp}}}\rangle_f}{S_f},\\
 \Delta^{Y}_{f,b}&=100\left\langle Y^{R_{\mathrm{resp}}}-Y^b\right\rangle_f,
\end{align}
where $Y$ is the final binary outcome. The family scales, fixed in advance, are $S_{\mathrm{Camera}}=0.316943$ and $S_{\mathrm{Sensor}}=0.200099$ in the same command units. They are constants, so $G$ is in units of $S_f$, not a percentage of each case's error. Main Figs.~\ref*{M-fig:teaser}b (against $N_T$) and~\ref*{M-fig:native-restoration}a (against all three) show these contrasts in command units, with each gain and interval endpoint multiplied by $S_f$.

\paragraph{Statistics: six simultaneous comparisons for each endpoint.}
Each endpoint has its own family of six comparisons: three comparators under two perturbation families. The prespecified procedure uses 10,000 bootstrap draws of whole tasks, keeping all paired observations of each drawn task. Within a family, the same task draws serve all comparators and both endpoints. The simultaneous 95\% intervals use percentile tails $0.05/(2\times6)$ and $1-0.05/(2\times6)$, with linear interpolation between ordered draws.

\subsection{Results: command error and task success}

\paragraph{Command error falls against the unedited run and both controls.}
All six command intervals are above zero (Table~\ref{tab:allstep-restoration-contrasts}). Against $N_T$, restoration lowers the task-equal command error from $0.4093$ to $0.3120$ under Camera and from $0.2091$ to $0.1987$ under Sensor Noise. In units of $S_f$, it reduces the error by $0.2908$ and $0.0337$ more than the clean-feature control, and by $0.2999$ and $0.0516$ more than the random control.

\paragraph{The reduction is broad across tasks.}
The table below summarizes the task-level command-error reductions behind Table~\ref{tab:allstep-restoration-contrasts}, as comparator minus $R_{\mathrm{resp}}$ in command units (the units of main Figs.~\ref*{M-fig:teaser}b and~\ref*{M-fig:native-restoration}a), with tasks weighted equally. Against $N_T$, the error falls in 25 of 27 Camera tasks and 24 of 30 Sensor Noise tasks. Medians lie below the means: the three largest task reductions account for 37\% (Camera) and 48\% (Sensor Noise) of the summed reduction. These are descriptive task statistics; the simultaneous intervals are those of Table~\ref{tab:allstep-restoration-contrasts}.

\begin{center}
\small
\setlength{\tabcolsep}{3pt}
\begin{tabular}{@{}lrrrrrr@{}}
\toprule
 & \multicolumn{3}{c}{Camera (27 tasks)} & \multicolumn{3}{c}{Sensor Noise (30 tasks)}\\
Comparator & Mean & Median & $>0$ & Mean & Median & $>0$\\
\midrule
$N_T$ & 0.0972 & 0.0557 & 25 & 0.0105 & 0.0060 & 24\\
$R_{\mathrm{clean}}$ & 0.0922 & 0.0538 & 23 & 0.0067 & 0.0041 & 23\\
$R_{\mathrm{rand}}$ & 0.0951 & 0.0552 & 25 & 0.0103 & 0.0057 & 24\\
\bottomrule
\end{tabular}
\end{center}


\paragraph{Under Camera, restoration moves the commands mainly toward the clean commands.}
Write the translation parts of the first five commands as $A^a\in\mathbb R^{15}$. For each setting, start state and noise draw, let $c=A^{N_C}-A^{N_T}$ be the direction toward the clean commands and $d^a=A^a-A^{N_T}$ the change the edit causes. Its signed component along $c$ and the size of the rest are
\begin{equation}
 d_\parallel^a=\frac{\langle d^a,c\rangle}{\|c\|_2},
 \qquad
 d_\perp^a=\left\|d^a-d_\parallel^a\frac{c}{\|c\|_2}\right\|_2.
 \label{eq:issued-command-direction}
\end{equation}
For each run, the command error satisfies $(e^a)^2=(\|c\|_2-d_\parallel^a)^2+(d_\perp^a)^2$: $d_\parallel$ is progress toward the clean commands and $d_\perp$ the change in other directions. Both are computed per record and then averaged like the command error.

All 1,296 records contain the five commands, and all 216 clean-target distances are positive. Table~\ref{tab:issued-command-direction} reports all six runs as a descriptive decomposition. Restoration has mean along/orthogonal components $0.10365/0.04660$ under Camera, so most of its change points toward the clean commands, and $0.01139/0.01229$ under Sensor Noise, where the change in other directions is as large as the progress toward the clean commands. Only the command-error contrasts carry intervals (Table~\ref{tab:allstep-restoration-contrasts}).

\begin{table*}[t]
\caption{\textbf{How each run changes the first five commands: progress toward the clean commands and change in other directions.} Descriptive task-equal means on all confirmation cases: 96 Camera cases (48 settings, 27 tasks) and 120 Sensor Noise cases (60 settings, 30 tasks), two noise draws per setting. Each quantity is computed per record from the translation parts of the first five commands, then averaged. $d_\parallel,d_\perp$ are defined in Eq.~\ref{eq:issued-command-direction}; $e$ is the command error. Command units; means without intervals. The geometric identity holds per record, not for the means.}

\centering\small
\begin{tabular}{lrrrrrr}\toprule
& \multicolumn{3}{c}{Camera} & \multicolumn{3}{c}{Sensor Noise}\\
\cmidrule(lr){2-4}\cmidrule(lr){5-7}
Run & Along $d_\parallel$ & Orthogonal $d_\perp$ & Error $e$ & Along $d_\parallel$ & Orthogonal $d_\perp$ & Error $e$\\\midrule
Clean ($N_C$) & 0.40928 & 0.00000 & 0.00000 & 0.20913 & 0.00000 & 0.00000\\
Perturbed, unedited ($N_T$) & 0.00000 & 0.00000 & 0.40928 & 0.00000 & 0.00000 & 0.20913\\
\midrule
Restoration ($R_{\mathrm{resp}}$) & 0.10365 & 0.04660 & 0.31205 & 0.01139 & 0.01229 & 0.19867\\
Clean-feature directions ($R_{\mathrm{clean}}$) & 0.00529 & 0.01001 & 0.40422 & 0.00388 & 0.00633 & 0.20541\\
Random directions ($R_{\mathrm{rand}}$) & 0.00221 & 0.00391 & 0.40710 & 0.00016 & 0.00246 & 0.20900\\
Whole L9 state ($R_{\mathrm{full}}$) & 0.11974 & 0.05176 & 0.29733 & 0.01671 & 0.01918 & 0.19433\\
\bottomrule\end{tabular}

\label{tab:issued-command-direction}
\end{table*}



\paragraph{The 16 directions reach most of what replacing the whole L9 state does.}
Replacing the whole L9 state lowers the mean command error to $0.2973$ (Camera) and $0.1943$ (Sensor Noise), only somewhat below restoration (Table~\ref{tab:allstep-restoration-arms}); the remaining error enters through the later blocks, which still read the perturbed image. These values are a reference, not one of the six comparisons. From Tables~\ref{tab:allstep-restoration-arms} and~\ref{tab:issued-command-direction}, replacing the whole state removes 27\% (Camera) and 7\% (Sensor Noise) of the command error of $N_T$, restoration 24\% and 5\%, so restoration reaches 86.9\% and 70.7\% of the whole-state reduction (from unrounded values, as in the main paper).

\paragraph{Edit size: the controls write less.}
The controls share restoration's number of directions, coefficient and schedule, but each copies only the clean--perturbed difference along its own directions; the shared directions, fitted to such differences, write five to seven times as much (Table~\ref{tab:allstep-restoration-arms}). Because a larger edit might by itself correct more, development compares the sets at matched edit size (Sec.~\ref{sec:dev-restoration-controls}), where restoration is better than clean-feature directions under Camera but worse under Sensor Noise. Restoration's edit size is 62\% and 71\% of the whole-state edit, so per unit of edit it removes 1.4 times as much command error as replacing the whole state under Camera and as much under Sensor Noise.

\begin{table*}[t]
\caption{\textbf{The six runs: command error, task success and edit size.} Descriptive task-equal means on the confirmation cases. $e$: command error (Eq.~\ref{eq:allstep-command-error}), zero for $N_C$ by definition. Edit size: actual change written into the L9 state (tokenwise $\ell_2$), summed over every query of the run's own episode, so it depends on both the writes and the episode length. Restoration and the controls share the number of directions, coefficient and schedule, yet write different amounts. $R_{\mathrm{full}}$ (whole L9 state) is a reference, not one of the six prespecified comparisons.}

\centering\small
\setlength{\tabcolsep}{9pt}
\begin{tabular}{@{}lrrrrrr@{}}
\toprule
 & \multicolumn{3}{c}{Camera} & \multicolumn{3}{c}{Sensor Noise}\\
\cmidrule(lr){2-4}\cmidrule(l){5-7}
Run & Error $e$ & Success (\%) & Edit size ($10^3$) & Error $e$ & Success (\%) & Edit size ($10^3$)\\
\midrule
$N_{\mathrm{C}}$ & 0.0000 & 98.15 & 0.00 & 0.0000 & 97.50 & 0.00\\
$N_{\mathrm{T}}$ & 0.4093 & 73.15 & 0.00 & 0.2091 & 83.33 & 0.00\\
$R_{\mathrm{resp}}$ & 0.3120 & 75.00 & 46.20 & 0.1987 & 86.67 & 51.47\\
$R_{\mathrm{clean}}$ & 0.4042 & 76.85 & 7.79 & 0.2054 & 83.33 & 7.74\\
$R_{\mathrm{rand}}$ & 0.4071 & 73.15 & 9.65 & 0.2090 & 85.00 & 9.64\\
$R_{\mathrm{full}}$ & 0.2973 & 80.56 & 74.20 & 0.1943 & 89.17 & 72.77\\
\bottomrule
\end{tabular}

\label{tab:allstep-restoration-arms}
\end{table*}


\paragraph{Task success is also higher.}
Observed success is higher with restoration: $75.00\%$ against $73.15\%$ under Camera and $86.67\%$ against $83.33\%$ under Sensor Noise. The clean runs succeed in $98.15\%$ and $97.50\%$ (Table~\ref{tab:allstep-restoration-arms}). The paired gains over $N_T$ are $+1.85$ and $+3.33$ percentage points. All six simultaneous success intervals include zero, the one against the clean-feature control under Sensor Noise with its lower end at zero (Table~\ref{tab:allstep-restoration-contrasts}). Under Camera, the clean-feature control has higher observed success than restoration.

\paragraph{Replacing the whole L9 state raises success.}
In a post hoc analysis with pointwise intervals, outside the six simultaneous comparisons, $R_{\mathrm{full}}$ minus $N_T$ is $+7.41\ [+0.93,+14.81]$ (Camera) and $+5.83\ [0.00,+12.50]$ (Sensor Noise): this layer matters for the outcome, and restoration delivers part of that gain.

\paragraph{Lost and gained successes in every run.}
The table counts, unweighted, the cases whose outcome differs from the unedited run $N_T$ with the same noise draw: failures that become successes (gained) and the reverse (lost). Gained successes grow with how much of the clean state is copied, while every edited run also loses some. The counts are descriptive; the simultaneous intervals are in Table~\ref{tab:allstep-restoration-contrasts}.

\begin{center}
\small
\setlength{\tabcolsep}{3pt}
\begin{tabular}{@{}lrrrrrr@{}}
\toprule
 & \multicolumn{3}{c}{Camera (96 cases)} & \multicolumn{3}{c}{Sensor Noise (120 cases)}\\
Run & Successes & Gained & Lost & Successes & Gained & Lost\\
\midrule
$N_T$ & 69 & -- & -- & 100 & -- & --\\
$R_{\mathrm{rand}}$ & 68 & 3 & 4 & 102 & 5 & 3\\
$R_{\mathrm{clean}}$ & 72 & 6 & 3 & 100 & 3 & 3\\
$R_{\mathrm{resp}}$ & 71 & 7 & 5 & 104 & 7 & 3\\
$R_{\mathrm{full}}$ & 76 & 10 & 3 & 107 & 10 & 3\\
$N_C$ & 94 & 26 & 1 & 117 & 19 & 2\\
\bottomrule
\end{tabular}
\end{center}


\subsection{Development checks: generation schedule and matched edit size}
\label{sec:dev-restoration-controls}

On the nine development tasks, restoring at every generation step matters. With edit size matched, restoration beats clean-feature directions under Camera but not under Sensor Noise; the confirmation runs use the unmatched edit ($\alpha=1$). These are descriptive means without intervals.

\paragraph{Saved states and plan endpoint.}
Camera covers eight settings on five tasks; Sensor Noise covers seventeen settings on nine tasks. Each setting has two noise draws and four saved query states, $q=0,2,5,10$, from the unedited perturbed run, giving 64 Camera and 136 Sensor Noise cases, all complete. Clean and perturbed observations share the saved physical state, every candidate generation starts from the same recorded policy noise, and each case keeps its assigned directions.

The endpoint is the first five positions of the generated translation plan, and its improvement is
\begin{equation}
 g^a=\frac{\|A^{N_T}-A^{N_C}\|_F-\|A^a-A^{N_C}\|_F}{S_f},
 \label{eq:dev-plan-gain}
\end{equation}
with the fixed family scales defined above. Means average the four saved queries and both noise draws within settings, then settings within tasks, with equal weight per task (Table~\ref{tab:dev-restoration-controls}). This endpoint is a plan at a saved state; the confirmation endpoint above is the first five commands of a complete rollout.

\begin{table}[t]
\caption{\textbf{Development: plan gains depend on the schedule and on the edit-size rule.} Descriptive task-equal means of Eq.~\ref{eq:dev-plan-gain} on all 64 Camera and 136 Sensor Noise cases: 8/17 settings on 5/9 tasks, two noise draws and saved queries $q=0,2,5,10$. First/all: edit at $j=0$ or at $j=0,\ldots,9$. Matched: Eq.~\ref{eq:dev-common-target} applied at each run's current state. $N_T$ has zero gain by definition.}

\centering\small
\begin{tabular}{lrr}\toprule
Rule & Camera / $S_f$ & Sensor Noise / $S_f$\\\midrule
Clean ($N_C$) & 1.20070 & 1.19909\\
Unedited ($N_T$) & 0.00000 & 0.00000\\
\midrule
Shared, first & 0.00475 & 0.00188\\
Shared, all & 0.15468 & 0.06687\\
Clean-feature, all & 0.02056 & 0.02137\\
Random, all & 0.00364 & 0.00118\\
Whole L9 state, all & 0.22249 & 0.12432\\
\midrule
Matched: shared & 0.02539 & 0.00975\\
Matched: clean-feature & 0.01763 & 0.01881\\
Matched: random & 0.00308 & 0.00044\\
\bottomrule\end{tabular}

\label{tab:dev-restoration-controls}
\end{table}


\paragraph{Restoring at every generation step matters.}
Both schedules use Eq.~\eqref{eq:allstep-restoration} with the shared directions. Moving from restoring only at $j=0$ to restoring at $j=0,\ldots,9$ raises the mean gain from $0.00475$ to $0.15468$ for Camera and from $0.00188$ to $0.06687$ for Sensor Noise. The schedule changes both the sequence of writes and the later generation states. The edit on original LIBERO, which acts only at the first step, is reported in Sec.~\ref{sec:supp-native-ablation}.

\paragraph{A common edit size at each position.}
The sets write equal amounts, so that they differ only in direction. For the shared, clean-feature and assigned random directions $B_k$, let $p_k=(H_T-H_C)B_kB_k^\top$ at one action position. Then
\begin{equation}
 s=\min_{k\in\{\mathrm{resp,clean,rand}\}}\|p_k\|_2,
 \qquad H'_k=H_T-s\frac{p_k}{\|p_k\|_2}.
 \label{eq:dev-common-target}
\end{equation}
A zero target leaves the state unchanged. The actual written changes, $H_{\rm after}-H_{\rm before}$, are recorded; their shared-minus-clean-feature norm differences take both signs. Under complete generation, each run reapplies this rule at its own current state.

\paragraph{Clean-feature edits change the local output more.}
Applying each candidate write at the same state of the unedited perturbed generation, we measure the change in the 32-dimensional normalized flow velocity, $\delta v_k=v_k-v_{N_T}$, separately for plan positions 0--4 and 5--9. Clean-feature edits give a larger $\|\delta v\|_2$ than shared-direction edits in every recorded mean (one per task, query, generation step and position group): $400/400$ for Camera ($5\times4\times10\times2$) and $720/720$ for Sensor Noise ($9\times4\times10\times2$). These grids repeat measurements of the same five and nine tasks. The clean local velocity was not recorded.

\paragraph{With edit size matched, the ranking differs between families.}
Under complete generation with matched edit size, the gains for shared/clean-feature directions are $0.02539/0.01763$ for Camera and $0.00975/0.01881$ for Sensor Noise. Restoration is thus better than clean-feature directions under Camera and worse under Sensor Noise, although clean-feature directions change the local output more in both: local output size is not plan correction.

\paragraph{Implementation notes.}
Projections and edit sizes are computed at higher precision than the policy's own execution, which writes the edited state at its native precision; writes that round to zero at that precision remain in the data.
